\begin{enumerate}
	\item $112\div 16=7$ et $140\div 16 =8,75$

	140 n'est pas divisible par 16, donc on ne peut pas constituer 16 sachets.
	\item \begin{align*}
	140&=2\times 70\\
	&=2\times 2\times 35\\
	&=2\times 2\times 5\times 7
	\end{align*}
	Donc la décomposition en produit de facteurs premiers de 140 est $2^2\times 5\times 7$.
	\item Grâce à leurs décompositions en produit de facteurs premiers, on remarque que 112 et 140 sont tous les deux divisibles par $2^2\times 7$ c'est-à-dire 28 ; et que c'est le plus grand diviseur commun de ces deux nombres.

	Le nombre maximum de sachets ayant la même composition qu'on peut constituer en utilisant tous les bonbons est donc de 28 sachets.

	$112\div 28 =4$ et $140\div 28=5$.

	Dans chaque sachet, il y aura 4 bonbons à la fraise et 5 bonbons au caramel.
\end{enumerate}
\needspace{10\baselineskip}
% saut de page si moins de 10 lignes disponibles
